% section: 問題を見たときの順番

\section{問題を解く前にすること}

漸化式を解くとき,最初から一般項を当てようとすると,式の形に
振り回されやすい.先に行うべきことは,次の項を作っている仕組みを
観察し,何を単純にすればよいかを決めることである.
ある程度慣れてきたらここにあることを無視して解答に走ってもらって構わないが,見慣れない問題の場合はこれらの手順を踏むと良い.

まずやるべきは初項がいくらかを見ることである.
例えば初項が$\displaystyle \frac{1}{\sqrt{6}-\sqrt{2}}$だとしよう.
わざわざこのような数字が設定されたということはきっと裏がある.
実際これは$\cos 15 \tcdegree $または$\sin 75 \tcdegree $の値である.
そこから考えて,この漸化式は三角関数型のものではないか,と予想するのである.

次にやることは$a_2$から$a_5$の値を求めてみることである.
最初の方の項を求めることで,法則性が見えてきたり,式がどのように移り変わっていくかが分かる.

それを終えたら次にやることは,式の概形を見ることである.